\subsection{Compact State and Event Timing Did Not Explain the Residual}

\begin{figure*}[t]
\centering
\includegraphics[width=\textwidth]{figures/Fig4_final.pdf}
\caption{Reduction boundaries. The tested $S(t)=(C(t),\dot{C}(t),R(t))$ first/second moment
closure and the state-matched event-local near-pre test did not provide stable
reductions of T1. Recent-history separation was visible in some observations
but did not become a universal sign/order rule. Propagation remained outside
the current confirmatory route.}
\label{fig:reduction_boundaries}
\end{figure*}

The compact-state reduction did not support a stable moment closure based on
$S(t)=(C(t),\dot{C}(t),R(t))$. The median incremental $R^2$ was approximately
$-6.33\times10^{-4}$ for the first-moment target and
$-6.21\times10^{-4}$ for the second-moment target, with low
positive-observation fractions. This result does not rule out all stochastic
or state-dependent descriptions. It shows that the tested compact-state
representation did not provide a stable first- or second-moment explanation of
T1.

The state-matched event-locality test produced a related boundary. True
transition timestamps did not show robust near-pre excess beyond
same-observation non-event frames matched in $S(t)$. The median
event-minus-control near-pre effect was approximately $-0.033$, and the
positive-observation fraction was about 0.42
(Fig.~\ref{fig:reduction_boundaries}). This does not mean that transitions
have no special dynamics; it means that the tested state-matched near-pre
aggregate did not provide a stable event-specific explanation.
